\section{Two linear stages}
\label{sec:composition}

A class bound for one stage is not closed under composition: even two
small-class maps can have many different endpoint rules.  We therefore
keep both complete operators.  Their exact composition forces a
\emph{common message} on a small rotation-invariant space.  Arithmetic
progressions then force this message space to grow.

\subsection{One message on several translates}

\begin{lemma}[Common-message growth]\label{lem:common-message}
\formal{CompositionProgression.lean}{commonMessage_finrank_le}
Let $A$ be a finite-dimensional vector space over any field, let
$R\colon A\to A$ be linear, and let $W\le A$ have dimension $D$.
Suppose that a linear map $L\colon A\to Z$ is injective on $W$ and
\[
 L(R^jx)=L(x)\qquad(x\in W,\ 0\le j<\ell),
\]
where $\ell\ge1$.  If $f=\dim\ker(R-I)$, then
\begin{equation}\label{eq:message-growth}
 \ell D\le \dim A+(\ell-1)f.
\end{equation}
Invertibility of $R$ is not required.
\end{lemma}
\begin{outline}
Put $V_t=W+RW+\cdots+R^{t-1}W$ and $V_0=0$.
For $t\le\ell$, the map $L$ vanishes on $(R-I)V_{t-1}$.
This image is disjoint from $W$, because $L$ is injective on $W$,
and their sum lies in $V_t$.  Rank--nullity therefore gives
\[
 \dim V_t\ge D+\dim V_{t-1}-f.
\]
Starting with $V_1=W$ and using $\dim V_\ell\le\dim A$ proves the
claim.  In fact telescoping gives
$V_t=W\oplus(R-I)V_{t-1}$, but the formal proof only needs inclusion.
\end{outline}

The correction $f$ matters.  Over $\F_2$, take a three-coordinate
rotation and $W=\operatorname{span}\{(1,1,1),e_0\}$.  Projection along
$\operatorname{span}\{(R-I)e_0\}$ identifies $W$ and $RW$ with the same
two-dimensional message.  Thus $\ell D=4=3+1$ for $\ell=2$; a bound
without the fixed-space term would be false.

For positive integers $M,\ell$, write $\mathsf W(M,\ell)$ for any
finite interval length such that every $M$-coloring of
$\{0,\ldots,\mathsf W(M,\ell)-1\}$ has a monochromatic ordinary
$\ell$-term arithmetic progression with positive step.
Existence is the finite van der Waerden theorem.  Our formal library
derives this version from Mathlib's finite Hales--Jewett theorem
\cite{HalesJewett1963}:
summing the symbols of a combinatorial line gives an ordinary progression
in a fixed bounded interval
(\lean{exists_progression_bound}).  We use no numerical estimate for
$\mathsf W$.

\begin{theorem}[Vanishing common-message rate]\label{thm:message-rate}
\formal{CompositionChannels.lean}{commonMessage_interval_bound}
Let $A=(F^q)^{\Z/d}$, with simultaneous cyclic rotation $T$, where
$q,d\ge1$ and $F$ is any field.  Let $W\le A$, and let
$B\colon W\to Z$ be linear and injective.  Suppose the $d$ positions
have $M$ nonempty classes and each class $J$ has a linear map
$L_J\colon A\to Z$ such that
\[
 L_J(T^cx)=B(x)\qquad(c\in J,\ x\in W).
\]
For $\ell\ge2$, put $N_0=\mathsf W(M,\ell)$.  If $d\ge N_0$, then
\begin{equation}\label{eq:message-rate}
 \dim W\le \frac{q(d+N_0-1)}{\ell}.
\end{equation}
\end{theorem}
\begin{outline}
Choose a monochromatic progression
$c_0,c_0+s,\ldots,c_0+(\ell-1)s$ among the first $N_0$ positions.
Here $s\ge1$ and $(\ell-1)s\le N_0-1$.
Apply \cref{lem:common-message} with $R=T^s$ and
$L=L_JT^{c_0}$.  A word fixed by $T^s$ is determined by its first
$s$ coordinates, by repeatedly reducing a natural-number index by $s$.
Thus the fixed space in $A$ has dimension at most $qs$ over every field.
No diagonalization is used.
Thus
\[
 \ell\dim W\le qd+(\ell-1)qs
 \le q(d+N_0-1).
\]
\end{outline}

For fixed $M$, this implies $\dim W/(qd)\to0$ uniformly in $F$, $q$,
the partition and the linear maps: choose $\ell$ first, then take
$d\ge\mathsf W(M,\ell)$, so the ratio is less than $2/\ell$.

Linearity here concerns the decoder, not merely the source.  To see the
distinction, identify binary words of length $d$ with $E=\F_{2^d}$ in a
normal basis, so rotation is Frobenius $\Phi$ \cite{Milne2022}.  Choose $\beta$ with
$d$ distinct conjugates.  The binary-linear two-channel space
$W=\{(x,\beta x):x\in E\}$ has dimension $d$, and its distinct rotated
views intersect only at zero.  A single total nonlinear decoder sends
$(u,\Phi^c(\beta)u)$, for $u\ne0$, to $\Phi^{-c}u$ and sends all other
points to zero.  It recovers the same original $x$ at every view with
rate $1/2$.  This is a source-level counterexample to removing decoder
linearity; it is not a two-stage transpose factorization.

\subsection{The complete two-stage reduction}

We now use the matrix model with $b$ coordinates per entry, allowing
arbitrary mixing of rows and lanes by every linear decoder.  Both stages
use the same column rotation.  Their partitions need not agree.

\begin{theorem}[Two linear stages]\label{thm:two-linear-stages}
\formal{CompositionLinear.lean}{twoHop_progression_gcd,twoHop_coprime_progression_dimension_le,twoHop_dimension_le,exists_twoHop_cutoff}
Let $K_1,K_2$ be linear class-uniform maps on $F^{d\times d\times b}$,
with $M_1,M_2$ nonempty classes, and suppose $K_2K_1=\tau$ on the
whole space.  If a first-stage class contains
$c_0,c_0+s,\ldots,c_0+(\ell-1)s$, with $\ell\ge1$, integer $s\ge0$
and indices read modulo $d$, then
\begin{equation}\label{eq:two-stage-ap}
 \ell d\le M_2d+(\ell-1)M_2\gcd(d,s).
\end{equation}
In particular, a progression of length $M_2+1$ with $\gcd(d,s)=1$
forces $d\le M_2^2$. Without the coprimality premise,
\begin{equation}\label{eq:two-stage-cutoff}
 d\le M_2\bigl(\mathsf W(M_1,M_2+1)-1\bigr).
\end{equation}
Thus fixed bounds on both class counts exclude exact two-stage
transposition at all sufficiently large $d$, over every field and
uniformly over every positive width $b$.
\end{theorem}
\begin{outline}
Both square factors are invertible, and $K_1^{-1}=\tau K_2$.
Let $S_0=\operatorname{Fix}(T)$, of dimension $bd$, and let $P_2\le F^d$
be the span of the second-class indicator vectors.  Their disjoint
nonempty supports give $\dim P_2=M_2$ over every field.
For $X\in S_0$, all views $T^cX$ coincide.  Hence $K_2X$ is constant
on each second receiver class.  It follows that
\[
 W:=K_1^{-1}S_0\ \le\
 A:=P_{2,\mathrm{row}}\otimes F^d_{\mathrm{column}}\otimes F^b,
 \qquad \dim W=bd,
 \qquad \dim A=bM_2d.
\]
This is an actual consequence of the global product identity, not an
independently chosen source.  The ambient $A$ is $T$-invariant and consists
of $bM_2$ cyclic channels.

The common column of $K_1X$, for $X\in W$, defines one injective linear
message $B(X)$.  At every receiver $c$ in a first class $J$,
$f_{1,J}(T^cX)=B(X)$.  Apply \cref{lem:common-message} on $A$ to the
specified progression. A word fixed by the $s$-step rotation is constant
on its $\gcd(d,s)$ position orbits. Restriction to one representative per
orbit is injective, so the fixed-space term is at most
$bM_2\gcd(d,s)$. The formal proof uses the actual translation-orbit quotient.
Canceling $b>0$ gives \cref{eq:two-stage-ap}; putting $\ell=M_2+1$ and
$\gcd(d,s)=1$ gives $d\le M_2^2$.

Put $N_0=\mathsf W(M_1,M_2+1)$.  If $d>M_2(N_0-1)$, then
$d\ge N_0$.  A monochromatic $(M_2+1)$-term progression in the first
$N_0$ positions has positive step $s$ with $M_2s\le N_0-1$.
Since $\gcd(d,s)\le s$ for this positive step,
\eqref{eq:two-stage-ap} would imply
\[
 d\le M_2^2s\le M_2(N_0-1),
\]
a contradiction.  This proves the finite cutoff.  Using fixed upper
bounds for the numbers of colors gives one cutoff for all the stated
parameters.
\end{outline}

The orbit correction can exclude particular masks before the general
Ramsey threshold applies. For example, with two second-stage classes,
a monochromatic three-term progression with step coprime to $d$ already
forces $d\le4$. The global two-class cutoff below remains $16$, since
Ramsey extraction does not require its progression step to be coprime.

The theorem concerns two complete maps on the same state space and
tensor side.  It neither multiplies endpoint class counts nor excludes
additional uniform postprocessing.  The familiar two partial-digit
swaps at $d=M_1M_2$ remain valid controls.  The bound does not prove
that their product condition is necessary for arbitrary coded factors.

\begin{corollary}[Two classes at each stage]\label{cor:two-binary-classes}
\formal{CompositionLinear.lean}{twoHop_two_classes_le_sixteen}
If each of the two linear stages uses at most two classes, exact
transposition requires $d\le16$, over every field and every positive width.
\end{corollary}
\begin{outline}
The finite statement $\mathsf W(2,3)\le9$ is checked by Lean's kernel
over all colorings of nine positions.  Its declaration is
\lean{binary_progression_nine_finite}; no minimality assertion is needed.
Substitute this bound in \cref{eq:two-stage-cutoff}.
\end{outline}

\subsection{One global rank error}

The preceding witness has dimension only $bd$.  To obtain more than a
row-scale obstruction, retain a whole band of frequencies and charge the
same global error only once.

\begin{lemma}[Recurrence growth]\label{lem:recurrence-growth}
\formal{CompositionRecurrence.lean}{recurrence_finrank_le}
Let $R\colon A\to A$ and $L\colon A\to Z$ be linear on a
finite-dimensional space, let $W\le A$ have dimension $D$, and let
$k,m\ge1$.  Let $p(z)=\sum_{i=0}^k a_i z^i$.
Suppose $(L,LR,\ldots,LR^{k-1})$ is injective on $W$ and
\[
 LR^jp(R)w=0\qquad(w\in W,\ j\ge0,\ j+2\le mk).
\]
Then, with $f=\dim\ker p(R)$,
\begin{equation}\label{eq:recurrence-growth}
 mD\le\dim A+(m-1)f.
\end{equation}
\end{lemma}
\begin{outline}
Set $V_t=\sum_{0\le j<t}R^jW$ and $G=(L,\ldots,LR^{k-1})$.
For $0\le t\le(m-1)k$, every component of $G$ vanishes on
$p(R)V_t$: its generator exponent is at most $t+k-2\le mk-2$.
Consequently $W\cap p(R)V_t=0$.
The degree bound gives $W+p(R)V_t\subseteq V_{t+k}$, so
\[
 \dim V_{t+k}\ge D+\dim V_t-f.
\]
Start with $\dim V_k\ge D$ and use $t=k,2k,\ldots,(m-1)k$.
Then $\dim V_{mk}\ge mD-(m-1)f$, proving the result.  When $m=k=1$
the displayed annihilation range is empty, as required.
\end{outline}

\begin{theorem}[Growing two-stage rank deficit]\label{thm:two-stage-rank}
\formal{CompositionRankField.lean}{twoHop_rank_interval_any_field,exists_twoHop_rank_threshold,twoHop_rank_superlinear}
Let the characteristic of $F$ not divide $d$, and let $K_1,K_2$ be
linear shared-decoder maps with $M_1,M_2$ classes on
$F^{d\times d\times b}$.  They need not be invertible.  Put
$r=\rank(\tau-K_2K_1)$.
For $k\ge1$, set
\[
 \ell=(2M_2+1)k-1,\qquad N_0=\mathsf W(M_1,\ell).
\]
If $d\ge N_0$, then
\begin{equation}\label{eq:rank-deficit}
 r\ge \frac{bk(d-N_0+1)}2.
\end{equation}
For fixed class bounds, $r/(bd)\to\infty$ uniformly in the field,
width, partitions and maps within this characteristic regime.
\end{theorem}
\begin{outline}
Extend scalars to a cyclotomic field with a primitive $d$th root
$\zeta$.  The complete maps and their receiver classes extend, and the
error rank does not increase: its extension factors through the scalar
extension of its original range.  It therefore suffices to prove the bound
there.  The characteristic hypothesis makes the roots distinct.  Let $C=\ker(\tau-K_2K_1)$ and $Y=K_1C$.
The restriction of $K_1$ to $C$ is injective, since $K_1x=0$ there
implies $\tau x=0$.  Thus $\dim Y\ge bd^2-r$, with $r$ denoting the
original error rank.

Let $S$ be the band of $k$ consecutive frequencies, consisting of
\[
 y_{:,c,:}=\sum_{\nu=0}^{k-1}\zeta^{\nu c}a_\nu.
\]
Its dimension is $bdk$, because $d\ge N_0\ge\ell\ge k$.
Pull $Y\cap S$ back through the injective restriction of $K_1$, obtaining
$W\le C$ of dimension $D\ge bdk-r$.  This intersection charges the
single global error $r$ once.

Let $H\le E^d$ be spanned by
$c\mapsto1_{p_2(c)=J}\zeta^{\nu c}$, over second classes $J$ and
$0\le\nu<k$, and put $h=\dim H\le M_2k$.
On $S$, the second operator produces columns of the form
$\sum_\nu\zeta^{\nu c}a_{p_2(c),\nu}$.
Since $w=\tau K_2K_1w$ on $W$, we have
\[
 W\le A:=H_{\mathrm{row}}\otimes E^d_{\mathrm{column}}\otimes E^b,
 \qquad \dim A=bhd.
\]
This ambient is invariant under column rotation.

Choose a monochromatic $\ell$-term progression with step $s\ge1$
among the first $N_0$ positions.  On $A$, put
$R=T^s$ and $L=f_{1,J}T^{c_0}$.  For $w\in W$,
\[
 LR^jw=\sum_{\nu=0}^{k-1}\zeta^{\nu c_0}
             (\zeta^{\nu s})^j a_\nu\qquad(0\le j<\ell).
\]
The numbers $\lambda_\nu=\zeta^{\nu s}$ are distinct:
for $\nu<\nu'$,
$0<(\nu'-\nu)s\le(k-1)s<d$.
The first $k$ displayed observations recover the $a_\nu$ by a
Vandermonde inverse, so their stacked map is injective on $W$.
The polynomial $p(z)=\prod_{\nu<k}(z-\lambda_\nu)$ annihilates these
sequences wherever all $k+1$ successive observations occur.  Taking
$m=2M_2$, the available exponents are exactly $j+2\le mk$.

The monic recurrence $p(R)x=0$ determines each cyclic channel from
its first $ks$ positions: its leading term recovers position $n$ from
positions $n-ks,n-(k-1)s,\ldots,n-s$.  Thus
$\dim\ker p(R)\le bhks$ on $A$.  This coordinate proof works without
diagonalizing the cyclic operator.
Apply \cref{lem:recurrence-growth} and $D\ge bdk-r$:
\[
 m(bdk-r)\le bhd+(m-1)bhks
 \le bM_2kd+(m-1)bM_2k^2s.
\]
Consequently
\[
 r\ge\frac{bkd}2-\frac{b(2M_2-1)k^2s}2.
\]
Now $(\ell-1)s\le N_0-1$ and
$(2M_2-1)k\le\ell-1$.  The subtracted term is at most
$bk(N_0-1)/2$, proving \cref{eq:rank-deficit}.

For each fixed $k$, once $d\ge\max\{N_0,2(N_0-1)\}$ the bound is
$r\ge bkd/4$.  Choosing $k$ arbitrarily large before $d$ proves the
uniform divergence of $r/(bd)$.
\end{outline}

This excludes repairs of uniformly $O(bd)$ rank, including all odd
binary dimensions.  It does not give $r\ge\epsilon bd^2$ for any fixed
$\epsilon>0$.  The arithmetic-progression cutoff grows with $k$, and
the proof does not permit $k$ to be a fixed positive fraction of $d$.
Nor is rank automatically a physical time charge.  Additional stages,
nonlinear factors and paid native realization require separate arguments.
